\section{The \texttt{MCPL\_input} McStas Component}
Source-like component that reads neutron state parameters from an mcpl-file.

\subsection*{Identification}
\begin{itemize}
  \item \textbf{Author:} Erik B Knudsen, Gregory S Tucker and Peter Willendrup
  \item \textbf{Origin:} DTU Physics / European Spallation Source ERIC
  \item \textbf{Date:} Mar 2016, consolidated with MCPL\_input\_once Oct 2026
\end{itemize}

\subsection*{Description}
Source-like component that reads neutron state parameters from a binary mcpl-file.

MCPL is short for Monte Carlo Particle List, and is a format for sharing events between e.g. MCNP(X), Geant4 and McStas.

The --ncount given on the commandline is overwritten by (\#particles in the file) x repeat\_count. Particles that are not neutrons, or are outside [Emin, Emax], are counted but ABSORBed.

The rays form a stream that runs through the file repeat\_count times. As in KDSource, this stream is split into one contiguous chunk per MPI process: each process seeks once to its start and reads on, wrapping around to the start of the file (with preload=1 it loads only the part of the file its chunk covers). Every event is thus used repeat\_count times across the whole run, and its weight is divided by repeat\_count. All passes after the first are smeared using v\_smear, pos\_smear and dir\_smear (the first pass is too if always\_smear is set), so repeat\_count \textgreater{} 1 MUST be combined with at least one *\_smear option; if all are 0, repeat\_count is reset to 1.

With replicate\_across\_processes=0 (the default, as in MCPL\_input\_once) the file is shared between the MPI processes and --mpi=N does not change the number of rays. replicate\_across\_processes=1 restores the pre-Oct 2026 MPI behaviour: repeat\_count is rounded up to a multiple of the number of MPI processes, so that every process runs through the whole file at least once, and all but the first pass on the master are smeared (or, with all *\_smear=0, are identical copies).

Example: MCPL\_input(filename=voutput,verbose=1,repeat\_count=1,v\_smear=0.1,pos\_smear=0.001,dir\_smear=0.01)

\subsection*{Input parameters}
Parameters in \textbf{boldface} are required; the others are optional.

\begin{longtable}{p{0.22\textwidth}p{0.12\textwidth}p{0.46\textwidth}p{0.14\textwidth}}
\toprule
\textbf{Name} & \textbf{Unit} & \textbf{Description} & \textbf{Default} \\
\midrule
\endhead
filename & str & Name of neutron mcpl file to read & 0 \\
polarisationuse &   & If !=0 read polarisation vectors from file & 1 \\
verbose &   & Print debugging information for the first 10 neutrons read (on the master process) & 0 \\
Emin & meV & Lower energy bound. Particles found in the MCPL-file below the limit are skipped & 0 \\
Emax & meV & Upper energy bound. Particles found in the MCPL-file above the limit are skipped & FLT\_MAX \\
repeat\_count & 1 & Replay the contents of the MCPL file this number of times, dividing the event weights by the same number. MUST be combined with \_smear options! & 1 \\
v\_smear & 1 & When replaying events, make a flat MC choice within v\_smear*V around particle velocity V & 0 \\
pos\_smear & m & When replaying events, make a flat MC choice of position within pos\_smear around particle starting position & 0 \\
dir\_smear & deg & When replaying events, make an MC choice of direction within dir\_smear around particle direction & 0 \\
preload &   & Load particles during INITIALIZE. On GPU preload is forced & 0 \\
always\_smear &   & If !=0, also smear the events on their first pass through the file & 0 \\
replicate\_across\_processes &   & If !=0, with MPI round repeat\_count up to a multiple of the number of processes, so that each process runs through the whole file & 0 \\
\bottomrule
\end{longtable}

\subsection*{Links}
\begin{itemize}
  \item Component source code found in file \texttt{MCPL\_input.comp}.
\end{itemize}
\IfFileExists{misc/MCPL_input_static.tex}{\input{misc/MCPL_input_static.tex}}{}